\section{Proof of the transfer theorem}\label{sec:proof}
We first separate an exact causal identity from its analytic estimate. Throughout, conditional states are taken under the law generated by the observer being constructed.

\begin{lemma}[All-incoming calibration and continuation identity]\label{lem:telescope}
Suppose the current labels have probabilities $p_i$ and conditional predictive states $q_i$. At label $i$ choose an action attaining the minimum in \eqref{eq:g}. Let $W_t$ be the post-report state before pooling after call $t$, and let $\lambda_t=\law(W_t)$. Partition $K$ into finitely many Borel cells $C_{t,j}$, and programme the next label $j$ with
\[
 q_{t,j}=\frac{\int_{C_{t,j}}z\,\lambda_t(dz)}{\lambda_t(C_{t,j})}.
\]
An arbitrary state is assigned to an unused label. Then $q_{t,j}$ is the actual predictive state conditional on label $j$. If the terminal decision minimizes $\ell(q_{n,j},u)$, its risk $R$ satisfies
\begin{equation}\label{eq:telescope}
 R-B_n^*=\sum_{t=1}^n\mathcal J_{V_t}(\lambda_t,\mathcal P_t).
\end{equation}
Each law $\lambda_t$ has density at most $\rho$ under Definition~\ref{def:hypotheses}(ii).
\end{lemma}
\begin{proof}
For any executable future test with affine expectation $f$, disintegration and the realization property give its expectation conditional on a retained label $i$, action $a_i$ and report $y$ as $f(F_{t-1}(q_i,a_i,y))$. Integration over every incoming edge assigned to cell $C_{t,j}$ yields $f(q_{t,j})$. Separation by tests proves calibration, and induction starts with $q_0=s_0$. There is no need to recover a discarded history. The unpooled law is the mixture $\sum_i p_i\Lambda_{t-1}(q_i,a_i)$; hence it obeys the density bound. Bellman optimality gives
\[
 \E V_t(W_t)=\E V_{t-1}(Q_{t-1}),\qquad
 Q_t=\E[W_t\mid Z_t].
\]
Subtracting and summing proves \eqref{eq:telescope}. At the last cut the minimizing readout has risk $\E g(Q_n)$. The representatives are constants computed for this very machine, including all incoming labels, not for a different full-history policy.
\end{proof}

\begin{lemma}[Boundary-valid acquired-mass curvature estimate]\label{lem:curvature}
Let $K\subset\R^d$ be compact, convex, and have nonempty interior. If $v$ is concave and Euclidean $L$-Lipschitz on $K$, then $-v$ has a globally convex $L$-Lipschitz extension $f$. For every finite measure $\lambda$ on $K$ and grid side $h>0$,
\begin{equation}\label{eq:curvbound}
 \mathcal J_v(\lambda,\mathcal P_h)
 \le C_d h^{2-d}\sum_C\lambda(C)
 \kappa_f(\overline B_{3\sqrt d h}(z_C)).
\end{equation}
If $\lambda$ is a probability law with density at most $\rho$, then for every integer $L_0\ge1$ there is a Borel partition with at most $L_0$ cells such that
\begin{equation}\label{eq:gridrate}
 \mathcal J_v(\lambda,\mathcal P)\le C_{K,d}\rho L L_0^{-2/d}.
\end{equation}
Neither an interior support margin nor differentiability of $v$ is required.
\end{lemma}
\begin{proof}
Put $f_0=-v$. At each $x\in\operatorname{int}K$ select all supporting slopes of $f_0$. Their Euclidean norms are at most $L$: directional difference quotients can be taken in every direction at an interior point. Take the supremum, on all $\R^d$, of the resulting supporting affine functions. This supremum is finite, convex, $L$-Lipschitz and agrees with $f_0$ on $\operatorname{int}K$; continuity gives equality on $K$. In particular, no assertion about arbitrary Lipschitz extensions preserving convexity is being used.

First suppose the extension is smooth. For a cell center $z$, subtract its tangent plane and call the resulting convex function $w$. Thus $w\ge0$ and $w(z)=0$. Put $R=2\sqrt d h$. For any $x$ in the cell, $B_{R/2}(x)\subset B_R(z)$. Writing $\operatorname{av}$ for the normalized spherical average, the submean inequality for $w$, radial monotonicity, and the divergence theorem give
\[
 w(x)\le \frac{C_d}{R^d}\int_{B_R(z)}w
 \le C_d\operatorname{av}_{\partial B_R(z)}w
 \le C_d R\operatorname{av}_{\partial B_R(z)}\partial_r w
 =C_dR^{2-d}\int_{B_R(z)}\Delta f.
\]
Here radial monotonicity gives $\int_{B_R}w\le (R/d)\int_{\partial B_R}w$.
The radial inequality is $w(z+R\omega)\le R\partial_r w(z+R\omega)$. This argument also applies in dimension one, with the boundary integral interpreted as the sum over two endpoints. If $b_C$ is the conditional barycenter, the affine tangent terms cancel and
\[
 \int_Cf\,d\lambda-\lambda(C)f(b_C)
 =\int_Cw\,d\lambda-\lambda(C)w(b_C)
 \le\lambda(C)\sup_Cw.
\]
Mollify $f$ on $\R^d$. The convex mollifications converge uniformly on bounded sets, and their positive Laplacian measures converge weakly to $\kappa_f$. The limsup of the mass of the radius-$R$ ball is bounded by the mass of the closed ball of radius $3\sqrt d h$. Summing the finitely many relevant cells proves \eqref{eq:curvbound}.

A Lipschitz convex function has uniformly bounded Laplacian mass on any fixed bounded region: choose a smooth nonnegative cutoff $\chi$ that equals one there and use
\[
 \int\chi\,d\kappa_f=-\int\nabla\chi\cdot\nabla f
 \le L\int|\nabla\chi|.
\]
For $h$ bounded above, the enlarged cell balls have bounded overlap, and $\lambda(C)\le\rho h^d$. Equation~\eqref{eq:curvbound} therefore gives $C_{K,d}\rho Lh^2$. Enclose $K$ in a fixed cube and subdivide each coordinate into $\lfloor L_0^{1/d}\rfloor$ intervals. There are at most $L_0$ cells. Since $\lfloor L_0^{1/d}\rfloor\ge L_0^{1/d}/2$, one has $h\le C_K L_0^{-1/d}$. This proves \eqref{eq:gridrate}, including $L_0=1$, by changing the fixed constant. The extension is used only in this proof, not as an observer state.
\end{proof}

\begin{lemma}[Continuation sensitivity and integer allocation]\label{lem:allocation}
With $a_t$ as in \eqref{eq:weights}, $\Lip(V_t)\le a_t$ in the chosen norm. If $m\ge1$ and $p=2/d$, there are positive integers $L_1,\ldots,L_n$ with $\sum_tL_t\le n+m-1$ and
\begin{equation}\label{eq:alloc}
 \sum_ta_tL_t^{-p}\le 2^p\left(\sum_ta_t^{1/(1+p)}\right)^{1+p}m^{-p}.
\end{equation}
Under \eqref{eq:stability}, $\sum_ta_t^{1/(1+p)}$ is bounded independently of $n$.
\end{lemma}
\begin{proof}
The Wasserstein inequality and the minimum over a common finite action set show $\Lip(V_t)\le\beta_t\Lip(V_{t+1})$; induction proves the first assertion. Concavity of $V_t$ follows independently by taking the infimum of the affine risks of all fixed continuation decision trees. Lipschitz continuity and a finite minimum give a Borel minimizing action.

Put $w_t=a_t^{1/(1+p)}$, $A=\sum_tw_t$, and choose
\[
 L_t=1+\left\lfloor\frac{(m-1)w_t}{A}\right\rfloor.
\]
When $m=1$, $\sum_ta_t\le A^{1+p}$. When $m\ge2$, $L_t\ge(m-1)w_t/A$, so the left side of \eqref{eq:alloc} is at most $A^{1+p}(m-1)^{-p}\le2^p A^{1+p}m^{-p}$. This counts one real label even at a very insensitive phase.

If $k=n-t$, decomposition into complete consecutive blocks gives
\[
 a_t\le L_g\max(1,B^{b-1})\theta^{\lfloor k/b\rfloor}.
\]
Summation of the resulting geometric series proves the uniform bound. No single-report contraction is asserted.
\end{proof}

\begin{lemma}[Clock support converse]\label{lem:clock}
Under Definition~\ref{def:hypotheses}(i), an autonomous observer stopping after exactly $n$ calls has pairwise disjoint positive-probability label supports at cuts $0,1,\ldots,n$. In particular, it has at least $n$ nonterminal labels and at most $M-n$ possible terminal labels. The claim allows fresh independent transition randomization.
\end{lemma}
\begin{proof}
Fix a label $z$ appearing at two different cuts $s<t$. From that same label, construct a reference observer trace by drawing each next report from $m^a$, where $a$ is the action selected by its current label, and by using its actual randomized label-transition kernel. For every finite number of future calls, the actual and reference trace measures are mutually absolutely continuous. Indeed, conditioning successively on each reference trace multiplies by strictly positive finite report densities almost everywhere. Mixtures over an unobserved incoming physical state or past history retain this equivalence. An unused value on a null report set is immaterial.

Starting at cut $t$ and label $z$, the machine stops after $n-t$ calls almost surely. The same event has reference probability one, and hence actual probability one starting at cut $s$ with label $z$. It would then stop before $n$, a contradiction. The argument also covers $t=n$, since stopping is a function of the label. Each of the $n+1$ cuts has nonempty support, proving the counts. Conversely, a deterministic chain of $n+1$ labels with any fixed legal sensing actions realizes the deadline. Its terminal forecast may be poor, but the class is nonempty.
\end{proof}

\begin{lemma}[Task curvature and actual-law converse]\label{lem:lower}
For every full-history policy, including random policies, the actual terminal predictive law has density at most $\rho$. If its terminal decision uses at most $L$ labels, then its risk is at least
\[
 B_n^*+c_{d,\alpha}\rho^{-2/d}L^{-2/d}.
\]
\end{lemma}
\begin{proof}
Condition immediately before the last action and report. For every selected incoming state and action the conditional post-report law has density at most $\rho$; mixing preserves that bound. Full-history randomization can be included among the independent observed variables for this argument. Its optimal full-history risk is at least $B_n^*$ by conditioning on independent seeds and minimizing over deterministic policies.

Let $S$ be the resulting full-history terminal state, and $Q=\E[S\mid Z]$ for the terminal label. Affinity of the conditional loss implies that the best label readout has risk $\E g(Q)$. Strong concavity gives
\[
 \E g(Q)-\E g(S)\ge\frac{\alpha}{2}\E|S-Q|^2.
\]
For any $L$ representatives, their radius-$r$ balls have total probability at most $\rho L v_dr^d$. Taking $r=(2\rho L v_d)^{-1/d}$ leaves probability at least $1/2$ outside their union, so the squared distortion is at least $r^2/2$. This applies also to randomized encoding, since $Q$ still takes at most $L$ values. No density lower bound or formal reachable-set packing is used.
\end{proof}

\begin{proof}[Proof of Theorem~\ref{thm:main}]
For the checkpoint upper, execute a Borel Bellman-optimal full-history policy and partition its actual terminal law using Lemma~\ref{lem:curvature} with $v=g$. The readout at each conditional barycenter gives $CM^{-2/d}$. Lemma~\ref{lem:lower} gives the matching checkpoint lower.

For an autonomous upper, put $m=M-n$ and take the allocation in Lemma~\ref{lem:allocation}. Starting with the preparation, inductively construct the calibrated observer in Lemma~\ref{lem:telescope}, using at most $L_t$ grid cells at phase $t$. Each grid is global on $K$ and therefore defines the next-label rule for every possible incoming label and report. The conditional barycenters belong to $K$ by convexity. Phase labels are disjoint, so the total state count is $1+\sum_tL_t\le M$. The actual density bound and the curvature estimate give
\[
 R-B_n^*\le C\sum_t a_tL_t^{-2/d}
 \le C A_n^{1/\gamma}m^{-2/d}.
\]
Norm-equivalence constants are absorbed in $C$. This proves the weighted assertion and, under block stability, the horizon-uniform upper. There is no inaccessible shadow state in its execution.

Finally Lemma~\ref{lem:clock} restricts every autonomous observer to at most $M-n$ terminal labels. Lemma~\ref{lem:lower} proves the lower in \eqref{eq:mainaut}. All optimization is over the declared task, preparation, action family and internal exact-deadline interface.
\end{proof}

\begin{remark}[Why the interface is part of the theorem]\label{rem:clock}
If the report at call $t$ includes an exact time stamp $t$, or an external terminal request is provided for free, positive trace equivalence across phases fails. A machine can then reuse labels across phases and need not pay $n$ distinct nonterminal states. Thus $M-n$ is not claimed for externally clocked filtering or for arbitrary stopping times. Likewise, without strong concavity a task may be insensitive to some acquired directions, and the stated full-dimensional lower need not hold. The mass--curvature upper remains meaningful in these cases.
\end{remark}
